\documentclass[11pt,a4paper]{article}
% Shared preamble for RuPhO English problem statements (match T1 2026 style).
% Use: \input{latex/rupho_en_preamble.tex} after \documentclass.
\usepackage[margin=1in]{geometry}
\usepackage{amsmath,amssymb}
\usepackage{graphicx}
\usepackage{float}
\usepackage{enumitem}
\usepackage{microtype}
\usepackage{caption}
\usepackage{tabularx}
\usepackage{hyperref}

\captionsetup{font=small,labelfont=bf,skip=6pt}

\setlist[itemize]{leftmargin=1.2cm}
\setlist[enumerate]{leftmargin=1.2cm}

% One outer box per subproblem: [id | statement | points] (no inner rules)
\newcommand{\problembox}[3]{%
  \par\addvspace{0.42em}%
  \noindent
  \setlength{\fboxsep}{7.5pt}%
  \setlength{\fboxrule}{0.95pt}%
  \fbox{%
    \begin{minipage}{\dimexpr\linewidth-2\fboxsep-2\fboxrule\relax}
    \begin{tabularx}{\linewidth}{@{}>{\raggedright\arraybackslash\bfseries}p{2.1em} @{\hspace{0.9em}} >{\raggedright\arraybackslash}X @{\hspace{0.9em}} >{\raggedleft\arraybackslash\bfseries}p{2.4em}@{}}
    #1 & #2 & #3
    \end{tabularx}
    \end{minipage}%
  }%
  \par\addvspace{0.32em}%
}


\title{\textbf{Satellite over the pole}}
\author{\normalsize Y13 (2013) --- English translation}
\date{}
\begin{document}
\maketitle

A satellite with an orbital period of $T = 4~\text{часа}$ is launched into orbit at the perigee point near the equator at low altitude. The earth's axis lies in the plane of this orbit.

\problembox{A1}{At what altitude $h$ and after what time $T_x$ after entering orbit can it fly over the pole for the first time?}{0}

The radius of the Earth is $R = 6 400~\text{км}$, the orbital period of a satellite flying in a low circular orbit is $T_0 = 85~\text{минут}$.

\vspace{1em}
\noindent\textit{Source:} \url{https://pho.rs/p/3996}

\end{document}